% BHU PYQ -- Auto-generated & error-corrected
\documentclass[11pt,a4paper]{article}

\usepackage[a4paper,top=2.5cm,bottom=2.5cm,left=2.8cm,right=2.8cm,headheight=14pt]{geometry}
\usepackage{amsmath,amssymb}
\usepackage[T1]{fontenc}
\usepackage[utf8]{inputenc}
\usepackage{lmodern}
\usepackage{microtype}
\usepackage{fancyhdr}
\usepackage{enumitem}
\usepackage{hyperref}
\usepackage{lastpage}
\usepackage{parskip}

\hypersetup{colorlinks=true,urlcolor=black,linkcolor=black,
  pdftitle={BSC-09A: Statistical Inference, Sampling and Design of Experiments -- BHU PYQ},pdfauthor={Banaras Hindu University}}

\pagestyle{fancy}\fancyhf{}
\renewcommand{\headrulewidth}{0.4pt}\renewcommand{\footrulewidth}{0.4pt}
\fancyhead[L]{\small   \textit{Banaras Hindu University}}
\fancyhead[R]{\small   \textit{BSC-09A: Statistical Inference, Sampling \& Design}}
\fancyfoot[C]{\small Page  \thepage\ of \pageref{LastPage}}


\newlist{parts}{enumerate}{2}
\setlist[parts,1]{label=(\alph*),leftmargin=*,itemsep=5pt,topsep=3pt}
\setlist[parts,2]{label=(\roman*),leftmargin=*,itemsep=3pt,topsep=2pt}

\newcommand{\pts}[1]{\hfill   \textit{[#1]}}


\begin{document}


\begin{center}
  {\large\bfseries BANARAS HINDU UNIVERSITY}\\[2pt]
  {
\normalsize B.Sc. (Hons.) Semester V Examination, 2013-14}\\[6pt]
  \rule{0.8\linewidth}{0.5pt}\\[6pt]
  {\large\bfseries Paper No. BSC-09A: Statistical Inference, Sampling and Design of Experiments}\\[4pt]
  {
\normalsize Time: 3 Hours\hfill Maximum Marks: 70}
\end{center}
\rule{\linewidth}{0.4pt}
\smallskip



\noindent   \textit{Instructions: Answer any   \textbf{five} questions. All questions carry equal marks. Symbols have their usual meanings.}
\bigskip

\medskip


\noindent   \textbf{Question 1.}

\begin{parts}
  \item Distinguish between sampling and complete enumeration. Explain the advantages of sampling over complete enumeration. \pts{7}
  \item Explain the concept of a sampling frame. How is it useful in selecting a sample? \pts{7}
\end{parts}

\medskip\rule{\linewidth}{0.2pt}

\medskip


\noindent   \textbf{Question 2.}

\begin{parts}
  \item Describe stratified random sampling. Write its advantages over simple random sampling and give the formula for the estimate of the population mean. \pts{7}
  \item Describe cluster sampling and two-stage sampling with examples. \pts{7}
\end{parts}

\medskip\rule{\linewidth}{0.2pt}

\medskip


\noindent   \textbf{Question 3.}

\begin{parts}
  \item What do you mean by design of experiment? Explain the principles of randomisation and replication. \pts{7}
  \item Describe the analysis of variance table for a Completely Randomised Design (CRD). \pts{7}
\end{parts}

\medskip\rule{\linewidth}{0.2pt}

\medskip


\noindent   \textbf{Question 4.}

\begin{parts}
  \item Describe the Latin Square Design (LSD) and derive its analysis of variance table. \pts{7}
  \item Describe the Randomised Block Design (RBD) along with its analysis of variance table. \pts{7}
\end{parts}

\medskip\rule{\linewidth}{0.2pt}

\medskip


\noindent   \textbf{Question 5.}

\begin{parts}
  \item Explain testing of hypothesis. In this context, describe the two kinds of errors (Type I and Type II) and the concept of level of significance. \pts{7}
  \item Explain two important uses of the $t$-test in statistics. \pts{7}
\end{parts}

\medskip\rule{\linewidth}{0.2pt}

\medskip


\noindent   \textbf{Question 6.}

\begin{parts}
  \item Describe the various uses of the $F$-test. \pts{7}
  \item Explain proportional allocation and optimum allocation in stratified sampling. Derive the formula for optimum allocation. \pts{7}
\end{parts}

\medskip\rule{\linewidth}{0.2pt}

\medskip


\noindent   \textbf{Question 7.}

\begin{parts}
  \item Give an example where an experiment is not properly designed and explain how proper design of experiments can remedy such deficiencies. \pts{7}
\end{parts}

\medskip\rule{\linewidth}{0.2pt}

\medskip


\noindent   \textbf{Question 8.} Write short notes on any   \textbf{two} of the following:\pts{14}

\begin{parts}
  \item Systematic sampling
  \item Factorial experiments
  \item $F$-test
  \item Multistage sampling
\end{parts}

\vfill
\rule{\linewidth}{0.4pt}

\begin{center}\small
  \href{https://bhu-exam.vercel.app/}{Click here to visit our website}\\[4pt]
  \href{https://me-aryan.vercel.app/}{Click here to visit our contributor}\\[4pt]
  \href{https://quashan-me.vercel.app/}{Click here to visit our contributor}
\end{center}

\end{document}